\subsection{不变元素}

定义 3. 设 \(\mathfrak{g}\) 是李代数，M 是 \(\mathfrak{g}\)-模。如果对所有 \(m \in \mathbf{M}\) 都有 \(\mathfrak{g}\)，则称元素 \(\mathfrak{g}\) 是不变的（相对于 M 上的 \(x_{\mathbf{M}}.m = 0\)-模结构，或相对于 \(x \in \mathfrak{g}\) 的相应表示）。

*设 G 是连通实李群，g 是其李代数，\(\theta\) 是 G 在有限维实向量空间 E 上的解析表示，\(\rho\) 是 g 在 E 上的相应表示。设 \(m \in E\)。元素 m 相对于 \(\rho\) 不变，当且仅当对所有 \(\theta(g).m = m\) 都有 \(g \in G\)。这说明使用“不变”一词是合理的。*

例 1. 设 M、N 是两个 g-模，且 \(P = \mathcal{L}_{K}(M, N)\)。由 (6)，P 的元素 f 不变的充要条件是 f 为从 g-模 M 到 g-模 N 的同态。特别地，若 M = N 且对所有 \(x_{M} = x_{N}\) 都有 \(x \in g\)，则 f 不变当且仅当 f 与各个 \(x_{M}\) 可交换。

例 2. 设 M 是有有限基的 K-模。若 M 具有 g-模结构，则 \(\mathcal{L}(\mathrm{M},\mathrm{M})\) 和 \(\mathbf{M}^*\otimes \mathbf{M}\) 都具有 g-模结构，而 \(\mathbf{M}^{*}\otimes \mathbf{M}\) 到 \(\mathcal{L}(\mathrm{M},\mathrm{M})\) 的典范映射是 g-模同构（命题 4）。由于 \(1\in \mathcal{L}(\mathrm{M},\mathrm{M})\) 显然是不变元（参见例 1），所以 \(u\) 中相应的元素 \(\mathbf{M}^{*}\otimes \mathbf{M}\) 也是不变元。若 \((e_i)_{1\leq i\leq n}\) 是 M 的基，\((e_{*}^{i})_{1\leqslant i\leqslant n}\) 是对偶基，则有 \(u = \sum_{i = 1}^{n}e_{i}^{*}\otimes e_{i}\)。

例 3。设 M 是 g-模。设 \(\beta\) 是 M 上的双线性形式，\(f\) 是 \(\mathcal{L}(\mathrm{M},\mathrm{M}^*)\) 中相应的元素。\(\beta\) 不变的充要条件是 \(f\) 为 g-模同态（命题 4 和例 1）。假设 K 是域且 \(\dim_{\mathbb{K}}\mathrm{M} < +\infty\)。M 上的非退化不变双线性形式 \(\beta\) 定义了从 g-模 M 到 g-模 \(\mathrm{M}^*\) 的同构，因而定义了从 g-模 \(\mathrm{M}\otimes \mathrm{M}\) 到 g-模 \(\mathrm{M}^*\otimes \mathrm{M}\) 的同构。因此，根据例 2，给定 \(\beta\) 便在 g-模中典范地定义了一个不变元 \(c\)，其所在空间为 \(\mathrm{M}\otimes \mathrm{M}\)，构造如下：设 \((e_i)_{1\leqslant i\leqslant n}\) 是 M 的基，\((e_i')_{1\leqslant i\leqslant n}\) 是满足 \(\beta (e_i,e_j') = \delta_{i,j}\) 的 M 的基；则 \(c = \sum_{i=1}^{n}e_i\otimes e_i'\)。

命题 5. 设 g 是李 K-代数，h 是 g 的理想，\(\rho\) 是 g 在 M 上的表示，\(\rho'\) 是 \(\rho\) 在 h 上的限制。则 M 中相对于 \(\rho'\) 不变的元素组成的集合 N 在 \(\rho(\mathfrak{g})\) 下稳定。

设 \(n\in \mathbf{N}\)，\(y\in \mathfrak{g}\)；对所有 \(x\in \mathfrak{h}\)，\([x,y]\in \mathfrak{h}\)，因而

\[
\rho (x) \rho (y) n = \rho ([ x, y ]) n + \rho (y) \rho (x) n = 0;
\]

所以 \(\rho(y)n\in\mathbf{N}\)。

命题 6. 设 M 是半单 g-模。则 M 的不变元子模 \(M_{0}\) 有唯一的在各个 \(x_{M}\) 下稳定的补模，即由 \(M_{1}\)（\(x_{M}.m\)）生成的子模 \(x \in g, m \in M\)。

设 \(\mathbf{M}'\) 是 \(\mathbf{M}\) 的一个在各个 \(x_{\mathrm{M}}\) 下稳定的子模，并且是 \(\mathrm{M}_0\) 在 \(\mathbf{M}\) 中的补模。对所有 \(m \in \mathbf{M}\)，都有 \(m = m_0 + m'\)，其中 \(m_0 \in \mathrm{M}_0\)、\(m' \in \mathrm{M}'\)，因而 \(x_{\mathrm{M}}m = x_{\mathrm{M}}m' \in \mathrm{M}'\)。所以 \(\mathrm{M}_1 \subset \mathrm{M}'\)。设 \(\mathrm{M}_2\) 是 \(\mathbf{M}'\) 的一个在各个 \(x_{\mathrm{M}}\) 下稳定的子模，并且是 \(\mathrm{M}_1\) 在 \(\mathrm{M}'\) 中的补模。对所有 \(m \in \mathrm{M}_2\)，

\[
x _ {\mathrm{M}} m \in \mathrm{M} _ {2} \cap \mathrm{M} _ {1} = \{0 \}
\]

对所有 \(x \in \mathfrak{g}\) 成立，故 \(m \in \mathbf{M}_0\)，从而 \(m = 0\)。因此 \(\mathbf{M}_2 = \{0\}\)，这就证明了 \(\mathbf{M}_1 = \mathbf{M}'\)。

